\chapter{はじめに}\label{章：はじめに}
\begin{bookcolumns}
\section{本書で扱う内容と対象読者}\label{節：内容と対象読者}

漸化式を解くとき,\,各項からある定数を引いたり,\,二つの項をまとめたりする.
うまく量を選べば,\,複雑に見えた式が等比数列の式に変わる.
しかし,\,計算ができた後にも疑問は残る.
なぜその量なのか. 特性方程式の根は何を表しているのか.
実数の数列を調べているのに,\,なぜ複素数や三角関数が現れるのか.

本書では,\,これらの疑問を線形代数と複素平面から考える.
中心となるのは,\,\textbf{同じ更新を,\,追いやすい量や座標で記録し直す}という見方である.
項を直接追う代わりに,\,不動点からのずれ,\,複数項の組,\,二つの差の比などを追う.
変えるのは問題そのものではなく,\,変化を記録する方法である.

第2章では,\,複数の値を一つの状態にまとめ,\,行列で更新する.
毎回同じ倍率で変化する成分に分けられれば,\,漸化式は複数の等比型へ分かれる.
第3章では,\,複素数の積を回転と拡大・縮小として読む.
分数の式,\,振動する数列,\,三角関数の公式が,\,平面上の動きへつながる.

対象は,\,高校数学の基本的な式計算を学び,\,漸化式の解法の背景に関心を持つ読者である.
行列や固有値を学んだ経験は必要としない.
数列と漸化式の基本も本文で確認する.
三角関数の基本的な意味は前提にするが,\,本書で使う加法公式と複素数との対応は説明する.
微分積分は使わない.

本書は『漸化式の体系と解法』と題材を共有するが,\,単独で読めるように構成した.
一般項を求めることに加えて,\,周期や収束,\,変形の意味を理解することを目標とする.

\section{本書の読み方}\label{節：本書の読み方}

まず具体的な更新を計算し,\,その計算を説明するために新しい概念を導入する.
一般の式を得た後には,\,初めの例へ戻って,\,見方の変化を確かめる.
例えば第2章の2階漸化式は第3章でも登場し,\,そこで複素平面上の動きとして読み直す.
一度読んだ式に戻ることが,\,二つの章を結ぶ.

例題は,\,原則として【例題】【方針】【解答】【解法】の順に置く.
解答だけでなく,\,どの量を作れば更新が簡単になるかを考えてほしい.
式を得たら,\,初期条件と更新式の両方に合うことまで確かめる.

左罫線の枠には,\,その箇所で持ち帰ってほしい見方をまとめる.
枠の式だけを暗記する必要はない.
どの条件で使え,\,何を記録している式なのかを本文で確かめよう.

行列が初めてなら第2章を順に読み,\,その後に第3章へ進むとよい.
複素平面に関心があれば第3章から始めてもよい.
最後の節では第2章の行列を使うので,\,その時点で必要な箇所へ戻れるよう参照を付けている.

\section{記号と表記の約束}\label{節：記号と表記}

数列を$a_1,a_2,\ldots$と書き,\,第$n$項を$a_n$と書く.
特に断らない限り,\,添字$n$は正の整数である.
第0項から始める場合には,\,その都度$a_0$と更新の範囲を明記する.

一般の階数には$k$,\,更新した回数には$m$を使う.
例えば
\[
  a_{n+2}=s a_{n+1}+t a_n
\]
は直前の二項が必要な\textbf{2階の漸化式}であり,\,式に三つの項が並ぶので\textbf{3項間漸化式}ともいう.
一般の$k$階では,\,次の項の計算に過去$k$項を使う.
ここでの階数は式の形を指し,\,特殊な初期値ではもっと低い階数の式も満たすことがある.

太字の$\boldsymbol v$は,\,複数の数を並べたベクトルを表す.
行列には$A,M,P$などを使い,\,何も変えない行列を$I$と書く.
実数全体を$\mathbb R$,\,複素数全体を$\mathbb C$と書く.
虚数単位は$i$,\,複素数$z$の共役は$\overline z$,\,実部は$\operatorname{Re}z$とする.

文字$p,q,r,s$などは,\,特に断らない限り,\,項の番号によらない定数である.
各節で役割を定め直すので,\,異なる節の同じ文字が同じ値とは限らない.
分母が零になる値は,\,割り算をする前に確認する.
複素平面に無限遠点を加える話も登場するが,\,元の実数列の割り算まで定義できるようになったとは解釈しない.

また,\,本編でいう「特性方程式型」
\[
  a_{n+1}=pa_n+q
\]
の$\alpha=p\alpha+q$は,\,\textbf{変わらない値}を探す方程式である.
一方,\,線形代数の特性方程式は,\,\textbf{更新の倍率}を探す.
同じ呼び名が使われる場合も,\,何を求めているかを区別する.
\end{bookcolumns}
